% Vecteurs vitesse v3 et v4 de la balle qui tombe et variation de vitesse (corrigé de l'exercice 2, chapitre 13)
% Échelle : 1 cm pour 1 m/s
\begin{tikzpicture}[x=1cm,y=1cm,
  vit/.style={otNavy,line width=1pt,-{Stealth[length=2mm]}},
  dv/.style={otGreenDark,line width=1.2pt,-{Stealth[length=2mm]}}]
\draw[otMuted,line width=0.4pt] (-0.4,0) -- (2.2,0);
\draw[vit] (0,0) -- (0,-1.72) node[pcLegende,anchor=east,text=otNavy,pos=0.5] {$\vec{v}_3$ (\qty{1,72}{m.s^{-1}})};
\draw[vit] (0.9,0) -- (0.9,-2.22) node[pcLegende,anchor=west,text=otNavy,pos=0.35] {$\vec{v}_4$ (\qty{2,22}{m.s^{-1}})};
\draw[otMuted,line width=0.4pt,dash pattern=on 1.5pt off 1.2pt] (0,-1.72) -- (1.8,-1.72) (0.9,-2.22) -- (1.8,-2.22);
\draw[dv] (1.8,-1.72) -- (1.8,-2.22) node[pcLegende,anchor=west,text=otGreenDark,pos=0.5] {$\Delta\vec{v}_3=\vec{v}_4-\vec{v}_3$ : \qty{0,50}{m.s^{-1}}};
\draw[otInk,line width=0.9pt] (5.2,-0.4) -- (5.2,-1.4) node[pcLegende,anchor=west,pos=0.5] {\qty{1}{m.s^{-1}}};
\end{tikzpicture}
\par\smallskip\legende{Vecteurs vitesse verticaux, vers le bas, tracés depuis un même niveau : $\Delta\vec{v}_3$ est vertical, vers le bas.}
